\section{Word Ladder I}
\label{sec:graph-word-ladder-1}

\problemstatement{Given a $\textit{beginWord}$, an $\textit{endWord}$,
and a $\textit{wordList}$, find the length of the \textbf{shortest}
transformation sequence from $\textit{beginWord}$ to
$\textit{endWord}$, changing exactly one letter at a time, where every
intermediate word must exist in $\textit{wordList}$. Return $0$ if no
such sequence exists.}

\begin{patternbox}
\emph{``Shortest sequence of single-character transformations''} is
shortest-path-in-an-unweighted-graph -- exactly what BFS computes --
but the graph is entirely \textbf{implicit}: there is no adjacency
list. A word's neighbors are generated on the fly by trying every
possible single-character substitution at every position and checking
whether the result is a valid word in the list.
\end{patternbox}

\subsection*{Generating neighbors on the fly}

For a word of length $L$ using a 26-letter alphabet, there are
$26L$ candidate single-character substitutions to try at each step.
Each candidate is checked against a hash set of valid words (built
from $\textit{wordList}$) in $O(L)$ (for the hash and string
comparison) -- so generating all neighbors of one word costs
$O(26L^2)$, in place of looking up a precomputed adjacency list.

\subsection*{Algorithm}

\begin{enumerate}
  \item Build a hash set from $\textit{wordList}$ for $O(1)$-ish
        membership checks (and remove words from the set as they are
        used, which serves as the \textit{visited} marker).
  \item BFS from $\textit{beginWord}$, tracking the transformation
        count alongside each word in the queue.
  \item For each word popped, try every single-character substitution
        at every position; if the result is in the set, it's a valid
        next step -- remove it from the set (mark visited) and push it
        with count $+1$. If the result equals $\textit{endWord}$,
        return the count $+1$ immediately.
  \item If the queue empties without reaching $\textit{endWord}$,
        return $0$.
\end{enumerate}

\subsection*{C++ implementation}

\begin{lstlisting}
#include <bits/stdc++.h>
using namespace std;

int ladderLength(string beginWord, string endWord, vector<string>& wordList) {
    unordered_set<string> words(wordList.begin(), wordList.end());
    if (words.find(endWord) == words.end()) return 0;

    queue<pair<string,int>> q;
    q.push({beginWord, 1});
    words.erase(beginWord);

    while (!q.empty()) {
        auto [word, steps] = q.front();
        q.pop();

        if (word == endWord) return steps;

        for (int i = 0; i < (int)word.size(); i++) {
            char original = word[i];
            for (char c = 'a'; c <= 'z'; c++) {
                if (c == original) continue;
                word[i] = c;
                if (words.find(word) != words.end()) {
                    words.erase(word);
                    q.push({word, steps + 1});
                }
            }
            word[i] = original;
        }
    }
    return 0;
}

int main() {
    string beginWord = "hit", endWord = "cog";
    vector<string> wordList = {"hot", "dot", "dog", "lot", "log", "cog"};
    cout << ladderLength(beginWord, endWord, wordList) << endl; // 5
    return 0;
}
\end{lstlisting}

\subsection*{Dry run}

\dryrun{Take $\textit{beginWord}=\texttt{"hit"}$,
$\textit{endWord}=\texttt{"cog"}$,
$\textit{wordList}=\{\texttt{hot,dot,dog,lot,log,cog}\}$.}

Queue: $[(\texttt{hit},1)]$. Pop \texttt{hit}: trying all
substitutions, \texttt{hot} matches (set has it); push
$(\texttt{hot},2)$, erase \texttt{hot}. Pop \texttt{hot}: substitutions
\texttt{dot} and \texttt{lot} both match; push
$(\texttt{dot},3),(\texttt{lot},3)$. Pop \texttt{dot}: \texttt{dog}
matches; push $(\texttt{dog},4)$. Pop \texttt{lot}: \texttt{log}
matches; push $(\texttt{log},4)$. Pop \texttt{dog}: \texttt{cog}
matches; push $(\texttt{cog},5)$. Pop \texttt{log}: \texttt{cog}
already erased from the set (claimed by \texttt{dog}'s expansion),
nothing new. Pop \texttt{cog}: $\texttt{word}=\textit{endWord}$,
return $5$.

\begin{complexitybox}
\textbf{Time Complexity.} $O(N \cdot 26L^2)$, where $N$ is the number
of words and $L$ is word length -- each of up to $N$ words is popped
once, and generating + checking its $26L$ candidate substitutions costs
$O(26L^2)$ (the $L$ factor from string copying/hashing each candidate).
\textbf{Space Complexity.} $O(N \cdot L)$ for the word set and the
queue.
\end{complexitybox}

\begin{takeawaybox}
When a problem's ``graph'' has no explicit edge list and instead
defines adjacency through some transformation rule (here, single-character
substitution), generate neighbors on demand inside the BFS loop rather
than trying to materialize the whole graph up front -- the BFS
structure itself is completely unchanged from the explicit-adjacency-list
case.
\end{takeawaybox}
